\documentclass[12pt,a4paper]{article}

% ---- Encoding & Font ----
\usepackage{iftex}
\ifPDFTeX
  \usepackage[utf8]{inputenc}
  \usepackage[T1]{fontenc}
\else
  \usepackage{fontspec}  % XeTeX/LuaTeX (tectonic): Unicode nativo (·, —, ¿, ª, §)
\fi
\usepackage[spanish]{babel}
\usepackage{csquotes}

% ---- Page layout ----
\usepackage[top=2.5cm, bottom=2.5cm, left=2.5cm, right=2.5cm]{geometry}
\usepackage{setspace}
\onehalfspacing

% ---- Math ----
\usepackage{amsmath, amssymb, amsthm}
\usepackage{mathtools}
\usepackage{physics}
\usepackage{esint}  % \oiint

% ---- Graphics ----
\usepackage{graphicx}
\usepackage{tikz}
\usepackage{float}
\usepackage{caption}

% ---- Tables ----
\usepackage{array}
\usepackage{booktabs}
\usepackage{tabularx}
\usepackage{colortbl}

% ---- Colours ----
\usepackage{xcolor}
\definecolor{notebg}{HTML}{FFF3CD}
\definecolor{noteborder}{HTML}{FFC107}
\definecolor{boxred}{HTML}{F8D7DA}
\definecolor{boxredborder}{HTML}{F5C6CB}

% ---- Boxes ----
\usepackage{mdframed}
\usepackage{tcolorbox}
\tcbuselibrary{most}

\newtcolorbox{keybox}[1][]{
  colback=blue!5,
  colframe=blue!70!black,
  arc=4pt,
  boxrule=1pt,
  fonttitle=\bfseries,
  title={#1}
}

\newtcolorbox{warnbox}[1][]{
  colback=boxred,
  colframe=boxredborder,
  coltitle=black!75,
  arc=4pt,
  boxrule=1pt,
  fonttitle=\bfseries,
  title={#1}
}

\newtcolorbox{notebox}[1][]{
  colback=notebg,
  colframe=noteborder,
  coltitle=black!75,
  arc=4pt,
  boxrule=1pt,
  fonttitle=\bfseries,
  title={#1}
}

% ---- Hyperlinks & TOC ----
\usepackage[hidelinks]{hyperref}
\usepackage{bookmark}

% ---- Enumerate ----
\usepackage{enumitem}

% ---- Miscellaneous ----
\usepackage{icomma}
\usepackage{siunitx}
\usepackage{textcomp}
\usepackage[bottom]{footmisc}

% ---- Diagramas (gráficas y esquemas) ----
\usepackage{pgfplots}
\pgfplotsset{compat=1.18}
\usetikzlibrary{babel}  % babel español activa `>`; esto evita romper las flechas `->`
% courses/ElecMag2024/Clases/assets/tikzstyles.tex
% ---------------------------------------------------------------------------
% Estilos compartidos para los diagramas del curso Electromagnetismo (2024).
% Fuente única del "look" visual (colores, grosores, escalas) para que todas
% las figuras (TikZ / pgfplots / CircuiTikz) sean consistentes entre clases.
%
% Es una COPIA de courses/Fisica3-2015/Clases/assets/tikzstyles.tex: mismo
% dominio, misma paleta. Copia y no symlink para que cada curso sea
% autocontenido y la paleta de Física III pueda evolucionar aparte.
%
% Uso: la clase debe cargar antes `tikz` y `pgfplots` (y `circuitikz` si la
%      clase tiene circuitos), y luego % courses/ElecMag2024/Clases/assets/tikzstyles.tex
% ---------------------------------------------------------------------------
% Estilos compartidos para los diagramas del curso Electromagnetismo (2024).
% Fuente única del "look" visual (colores, grosores, escalas) para que todas
% las figuras (TikZ / pgfplots / CircuiTikz) sean consistentes entre clases.
%
% Es una COPIA de courses/Fisica3-2015/Clases/assets/tikzstyles.tex: mismo
% dominio, misma paleta. Copia y no symlink para que cada curso sea
% autocontenido y la paleta de Física III pueda evolucionar aparte.
%
% Uso: la clase debe cargar antes `tikz` y `pgfplots` (y `circuitikz` si la
%      clase tiene circuitos), y luego % courses/ElecMag2024/Clases/assets/tikzstyles.tex
% ---------------------------------------------------------------------------
% Estilos compartidos para los diagramas del curso Electromagnetismo (2024).
% Fuente única del "look" visual (colores, grosores, escalas) para que todas
% las figuras (TikZ / pgfplots / CircuiTikz) sean consistentes entre clases.
%
% Es una COPIA de courses/Fisica3-2015/Clases/assets/tikzstyles.tex: mismo
% dominio, misma paleta. Copia y no symlink para que cada curso sea
% autocontenido y la paleta de Física III pueda evolucionar aparte.
%
% Uso: la clase debe cargar antes `tikz` y `pgfplots` (y `circuitikz` si la
%      clase tiene circuitos), y luego \input{../assets/tikzstyles.tex}
% (compilar desde el directorio de la clase para que la ruta relativa resuelva).
% ---------------------------------------------------------------------------

% ---- Paleta (alineada con las cajas del preámbulo: keybox/notebox/warnbox) ----
\definecolor{figblue}{HTML}{1F4E79}   % azul principal (líneas, keybox)
\definecolor{figred}{HTML}{C0392B}    % rojo de énfasis (warnbox)
\definecolor{figamber}{HTML}{B8860B}  % ámbar (notebox)
\definecolor{figgray}{HTML}{555555}   % auxiliares, ejes, guías

% ---- CircuiTikz: escala y grosor de los componentes ----
% Guarda \ifdefined: la electrostática (Clases 1-14) no carga `circuitikz`, y
% sin la guarda el \input revienta con `Undefined control sequence \ctikzset`.
% Cuando el curso llegue a circuitos, el bloque se activa solo.
\ifdefined\ctikzset
  \ctikzset{
    bipoles/thickness=1,
    resistors/scale=0.9,
    capacitors/scale=0.9,
  }
\fi

% ---- TikZ: estilos reutilizables ----
% Nota: las puntas se declaran como `-{Latex[...]}` y no como `->`. La forma
% con llaves NO contiene un `>` literal, así que es inmune al `>` activo que
% introduce babel español (ver CLAUDE.md §7.3.3).
\usetikzlibrary{arrows.meta,calc,angles,quotes}

\tikzset{
  figline/.style={figblue, line width=0.9pt},
  figaux/.style={figgray, dashed, line width=0.6pt},
  % vectores
  figvec/.style={figblue, line width=0.9pt, -{Latex[length=2.2mm,width=1.6mm]}},
  figvecres/.style={figred, line width=1.3pt, -{Latex[length=2.6mm,width=1.9mm]}},
  figvecaux/.style={figgray, line width=0.7pt, -{Latex[length=1.8mm,width=1.3mm]}},
  figaxis/.style={figgray, line width=0.6pt, -{Latex[length=2mm,width=1.4mm]}},
  % cargas puntuales
  figcarga/.style={circle, inner sep=0pt, minimum size=5.4pt, draw=none},
  figpos/.style={figcarga, fill=figred},
  figneg/.style={figcarga, fill=figblue},
  figneutra/.style={figcarga, fill=figgray},
  % etiquetas
  figlbl/.style={font=\small},
  figlblaux/.style={font=\footnotesize, figgray},
}

% ---- pgfplots: estilo base de las gráficas ----
\pgfplotsset{
  emfig/.style={
    width=11cm, height=6.5cm,
    axis lines=left,
    xlabel style={font=\small},
    ylabel style={font=\small},
    tick label style={font=\footnotesize},
    every axis plot/.append style={line width=1pt},
    grid=major,
    grid style={figgray!20},
    legend cell align=left,
    legend style={font=\footnotesize, draw=figgray!40},
  },
}

% (compilar desde el directorio de la clase para que la ruta relativa resuelva).
% ---------------------------------------------------------------------------

% ---- Paleta (alineada con las cajas del preámbulo: keybox/notebox/warnbox) ----
\definecolor{figblue}{HTML}{1F4E79}   % azul principal (líneas, keybox)
\definecolor{figred}{HTML}{C0392B}    % rojo de énfasis (warnbox)
\definecolor{figamber}{HTML}{B8860B}  % ámbar (notebox)
\definecolor{figgray}{HTML}{555555}   % auxiliares, ejes, guías

% ---- CircuiTikz: escala y grosor de los componentes ----
% Guarda \ifdefined: la electrostática (Clases 1-14) no carga `circuitikz`, y
% sin la guarda el \input revienta con `Undefined control sequence \ctikzset`.
% Cuando el curso llegue a circuitos, el bloque se activa solo.
\ifdefined\ctikzset
  \ctikzset{
    bipoles/thickness=1,
    resistors/scale=0.9,
    capacitors/scale=0.9,
  }
\fi

% ---- TikZ: estilos reutilizables ----
% Nota: las puntas se declaran como `-{Latex[...]}` y no como `->`. La forma
% con llaves NO contiene un `>` literal, así que es inmune al `>` activo que
% introduce babel español (ver CLAUDE.md §7.3.3).
\usetikzlibrary{arrows.meta,calc,angles,quotes}

\tikzset{
  figline/.style={figblue, line width=0.9pt},
  figaux/.style={figgray, dashed, line width=0.6pt},
  % vectores
  figvec/.style={figblue, line width=0.9pt, -{Latex[length=2.2mm,width=1.6mm]}},
  figvecres/.style={figred, line width=1.3pt, -{Latex[length=2.6mm,width=1.9mm]}},
  figvecaux/.style={figgray, line width=0.7pt, -{Latex[length=1.8mm,width=1.3mm]}},
  figaxis/.style={figgray, line width=0.6pt, -{Latex[length=2mm,width=1.4mm]}},
  % cargas puntuales
  figcarga/.style={circle, inner sep=0pt, minimum size=5.4pt, draw=none},
  figpos/.style={figcarga, fill=figred},
  figneg/.style={figcarga, fill=figblue},
  figneutra/.style={figcarga, fill=figgray},
  % etiquetas
  figlbl/.style={font=\small},
  figlblaux/.style={font=\footnotesize, figgray},
}

% ---- pgfplots: estilo base de las gráficas ----
\pgfplotsset{
  emfig/.style={
    width=11cm, height=6.5cm,
    axis lines=left,
    xlabel style={font=\small},
    ylabel style={font=\small},
    tick label style={font=\footnotesize},
    every axis plot/.append style={line width=1pt},
    grid=major,
    grid style={figgray!20},
    legend cell align=left,
    legend style={font=\footnotesize, draw=figgray!40},
  },
}

% (compilar desde el directorio de la clase para que la ruta relativa resuelva).
% ---------------------------------------------------------------------------

% ---- Paleta (alineada con las cajas del preámbulo: keybox/notebox/warnbox) ----
\definecolor{figblue}{HTML}{1F4E79}   % azul principal (líneas, keybox)
\definecolor{figred}{HTML}{C0392B}    % rojo de énfasis (warnbox)
\definecolor{figamber}{HTML}{B8860B}  % ámbar (notebox)
\definecolor{figgray}{HTML}{555555}   % auxiliares, ejes, guías

% ---- CircuiTikz: escala y grosor de los componentes ----
% Guarda \ifdefined: la electrostática (Clases 1-14) no carga `circuitikz`, y
% sin la guarda el \input revienta con `Undefined control sequence \ctikzset`.
% Cuando el curso llegue a circuitos, el bloque se activa solo.
\ifdefined\ctikzset
  \ctikzset{
    bipoles/thickness=1,
    resistors/scale=0.9,
    capacitors/scale=0.9,
  }
\fi

% ---- TikZ: estilos reutilizables ----
% Nota: las puntas se declaran como `-{Latex[...]}` y no como `->`. La forma
% con llaves NO contiene un `>` literal, así que es inmune al `>` activo que
% introduce babel español (ver CLAUDE.md §7.3.3).
\usetikzlibrary{arrows.meta,calc,angles,quotes}

\tikzset{
  figline/.style={figblue, line width=0.9pt},
  figaux/.style={figgray, dashed, line width=0.6pt},
  % vectores
  figvec/.style={figblue, line width=0.9pt, -{Latex[length=2.2mm,width=1.6mm]}},
  figvecres/.style={figred, line width=1.3pt, -{Latex[length=2.6mm,width=1.9mm]}},
  figvecaux/.style={figgray, line width=0.7pt, -{Latex[length=1.8mm,width=1.3mm]}},
  figaxis/.style={figgray, line width=0.6pt, -{Latex[length=2mm,width=1.4mm]}},
  % cargas puntuales
  figcarga/.style={circle, inner sep=0pt, minimum size=5.4pt, draw=none},
  figpos/.style={figcarga, fill=figred},
  figneg/.style={figcarga, fill=figblue},
  figneutra/.style={figcarga, fill=figgray},
  % etiquetas
  figlbl/.style={font=\small},
  figlblaux/.style={font=\footnotesize, figgray},
}

% ---- pgfplots: estilo base de las gráficas ----
\pgfplotsset{
  emfig/.style={
    width=11cm, height=6.5cm,
    axis lines=left,
    xlabel style={font=\small},
    ylabel style={font=\small},
    tick label style={font=\footnotesize},
    every axis plot/.append style={line width=1pt},
    grid=major,
    grid style={figgray!20},
    legend cell align=left,
    legend style={font=\footnotesize, draw=figgray!40},
  },
}


% ---- Title ----
\title{\textbf{Clase 10: Condiciones de frontera entre dieléctricos y Laplace con dos medios} \\[0.3em]
        \large{El signo de la susceptibilidad, la esfera en un dieléctrico y el cilindro en campo uniforme}}
\author{Transcripto y expandido --- Curso Electromagnetismo (OpenFING)}
\date{}

\begin{document}

\maketitle
\thispagestyle{empty}
\newpage

\tableofcontents
\newpage

\section{Consecuencias del signo de la susceptibilidad}

La clase anterior definió susceptibilidad y permitividad para materiales
homogéneos, isótropos y \textbf{lineales}, entendiendo por lineal que $\chi$ no
depende del campo aplicado y por lo tanto la relación entre $\vec P$ y $\vec E$
---o entre $\vec D$ y $\vec E$--- es de proporcionalidad.

Falta una observación elemental que tiene consecuencias fuertes. Si hay un campo
externo, está asociado a cargas exteriores positivas y negativas. La parte
negativa de cada molécula tiende a acercarse a la carga \textbf{positiva}
externa, y la positiva a la negativa. El momento dipolar que se genera apunta
entonces en el mismo sentido que el campo aplicado, y como la polarización es el
momento dipolar por unidad de volumen:

\begin{keybox}[Tres cotas encadenadas]
\[
\vec P \ \text{tiene el mismo sentido que}\ \vec E
\quad\Longrightarrow\quad
\boxed{\ \chi \ge 0\ }
\]
y usando $\varepsilon = \varepsilon_0(1+\chi)$,
\[
\boxed{\ \varepsilon \ge \varepsilon_0
\qquad\text{y}\qquad
K = \frac{\varepsilon}{\varepsilon_0} \ge 1 \ }
\]
\end{keybox}

\begin{notebox}[Origen trivial, conclusión no trivial]
El argumento es que las cargas negativas son atraídas por las positivas. Pero la
conclusión no es menor: \textbf{la constante dieléctrica no puede ser cualquier
cosa}, siempre es mayor o igual que uno.
\end{notebox}

\subsection{Materiales anisótropos}

En algunos sólidos el material no es isótropo. La generalización es casi la
misma, pero la susceptibilidad deja de ser un escalar y pasa a ser una
\textbf{matriz} ---un tensor--- que a un campo aplicado le asocia una
polarización que \textbf{no tiene por qué ser colineal} con él. Ese caso no se
trata en el curso; se menciona sólo para que se sepa que existe.

\subsection{Valores típicos}

\begin{table}[H]
  \centering
  \begin{tabular}{lc}
    \toprule
    Material & $K$ \\
    \midrule
    Azufre                          & 4,0 \\
    Cuarzo                          & 4,3 \\
    Cloruro de sodio (sal de mesa)  & 6,1 \\
    Agua destilada (\SI{20}{\celsius}) & 80,1 \\
    Aire (presión normal)           & $\approx 1$ \\
    \bottomrule
  \end{tabular}
  \caption{Valores indicativos del tamaño del efecto.}
\end{table}

\begin{notebox}[Dos lecturas de la tabla]
Primero: \textbf{el aire y el vacío son lo mismo} a los efectos de este curso
---distinguirlos exigiría una precisión que nunca se va a usar---. Segundo: para
materiales nada exóticos el efecto es \textbf{significativo}, con factores de 6 o
10.
\end{notebox}

\begin{warnbox}[Agua destilada, no agua con sales]
El agua con sales disueltas \textbf{no} es un aislante sino un conductor. El
valor de la tabla corresponde a agua destilada, y depende algo de la temperatura.
\end{warnbox}

\section{Primer ejemplo: esfera cargada en un dieléctrico}

El ejemplo más simple posible, con simetría esférica y dependencia de una sola
coordenada. Una \textbf{esfera conductora} de radio $a$ con carga total $Q$,
rodeada de un medio dieléctrico de constante $K$ (homogéneo, isótropo y lineal
--- todo lo que la frase «constante $K$» implica).

\begin{figure}[H]
  \centering
  % Clase 10 — esfera conductora cargada rodeada de dielectrico.
% El punto delicado del ejemplo, y el motivo de la figura, es la NORMAL: el
% material es el de afuera, asi que la normal saliente del material apunta hacia
% el centro, y de ahi el signo menos de sigma_p.
% CUIDADO CON LOS DOS VERSORES: e_r y n = -e_r nacen del mismo punto y en la
% misma recta, asi que dibujarlos en el mismo angulo hace que los rotulos se
% pisen (paso, y era ilegible). Van en angulos SEPARADOS, cada uno con su
% rotulo lejos del otro. La colinealidad no es el contenido aca --- el contenido
% es que apuntan al reves --- asi que separarlos no miente.
% La prosa explicativa va en el \caption, no adentro: adentro competia por el
% mismo espacio que los vectores.
\begin{tikzpicture}[scale=1]

  % dielectrico (corona)
  \fill[figamber!10] (0,0) circle (3.4);
  \draw[figamber, line width=0.8pt] (0,0) circle (3.4);
  \node[figlbl, figamber] at (-1.75,2.75) {diel\'ectrico $K$};

  % conductor
  \fill[figgray!25] (0,0) circle (1.1);
  \draw[figline, line width=1.0pt] (0,0) circle (1.1);
  \node[figlbl] at (0,0.28) {$Q$};
  \draw[figvecaux] (0,0) -- (160:1.1);
  \node[figlblaux, fill=figgray!25, inner sep=1pt] at (160:0.6) {$a$};

  % superficie gaussiana
  \draw[figline, dashed, line width=0.9pt, figblue] (0,0) circle (2.3);
  \node[figlbl, figblue, fill=figamber!10, inner sep=1.5pt] at (118:2.62) {$S$};
  \draw[figvecaux] (0,0) -- (300:2.3);
  \node[figlblaux, fill=figamber!10, inner sep=1pt] at (300:1.62) {$r$};

  % campo radial: solo donde no molesta a los versores
  \foreach \a in {90,140,190,270}{
    \draw[figvec] (\a:1.25) -- (\a:2.05);}
  \node[figlbl, figblue, fill=figamber!10, inner sep=1.5pt] at (163:2.95)
    {$\vec E,\ \vec D,\ \vec P$};

  % dipolos orientados, en un sector libre
  \foreach \a in {205,220,235}{
    \begin{scope}[shift={(\a:1.62)}, rotate=\a]
      \draw[figline, line width=0.7pt] (-0.16,0) -- (0.16,0);
      \node[figneg, minimum size=4pt] at (-0.16,0) {};
      \node[figpos, minimum size=4pt] at (0.16,0) {};
    \end{scope}}

  % versor radial: angulo 55
  \draw[figvecaux] (55:1.15) -- (55:2.0);
  \node[figlblaux, figgray, fill=figamber!10, inner sep=1.5pt] at (55:2.3)
    {$\hat e_r$};

  % normal saliente DEL MATERIAL: angulo 15, apunta al centro
  \draw[figvecres] (15:1.1) -- (15:0.42);
  \node[figlbl, figred, fill=figamber!10, inner sep=1.5pt] at (15:2.35)
    {$\hat n=-\hat e_r$};
  \draw[figaux] (15:1.95) -- (15:1.2);

  % sigma_p en la superficie del conductor
  \node[figlblaux, figblue, fill=figgray!25, inner sep=1.5pt] at (0,-0.72)
    {$\sigma_p<0$};

\end{tikzpicture}

  \caption{El material dieléctrico es el que está \emph{afuera}, así que la
  normal que sale de él apunta al centro: $\hat n = -\hat e_r$. De ahí que
  $\sigma_p<0$, y que la carga total que ve $\vec E$ sea $Q/K$ en vez de $Q$.}
\end{figure}

\subsection{Campo y polarización}

Se toma una \textbf{superficie gaussiana esférica} $S$ de radio $r > a$. Dentro
del conductor no hay campo, así que la región interesante es la exterior. Por
simetría esférica, $\vec E$, $\vec P$ y $\vec D$ son todos radiales y paralelos
entre sí ---el material es lineal e isótropo---, y sus módulos son constantes
sobre $S$. Aplicando Gauss para dieléctricos:
\[
\oiint_S \vec D \cdot \hat n \, dS = D \cdot 4\pi r^2 = Q
\quad\Longrightarrow\quad
\boxed{\ D = \frac{Q}{4\pi r^2}\ }
\]
donde $Q$ es la \textbf{carga libre}, que en este problema es toda la carga del
conductor: la que uno deposita y a la que tiene acceso.

Ahora hace falta volver a $\vec E$, y para eso se usa la relación constitutiva
del medio lineal, $\vec D = \varepsilon \vec E = \varepsilon_0 K \vec E$:

\begin{keybox}[Campo en el dieléctrico]
\[
\boxed{\ E = \frac{Q}{4\pi\varepsilon_0 K r^2}\ }
\]
menor que en el vacío, exactamente por un factor $K$.
\end{keybox}

\begin{notebox}[Por qué hacía falta la relación constitutiva]
Sin ella se podía calcular $\vec D$ por simetría, pero no habría servido de
mucho: la magnitud que ejerce la fuerza electrostática es $\vec E$, no $\vec D$.
\end{notebox}

Para la polarización se usa la otra forma de la relación,
$\vec D = \varepsilon_0 \vec E + \vec P$:
\[
\vec P = \vec D - \varepsilon_0 \vec E
= \varepsilon_0 (K-1) \vec E
\quad\Longrightarrow\quad
\boxed{\ P = \frac{(K-1)\,Q}{4\pi K r^2}\ }
\]

\begin{notebox}[Coherencia]
Si $K = 1$ estamos en el vacío y $P = 0$, como debe ser. Y como $K > 1$,
$\vec P$ tiene el mismo sentido que $\vec E$ ---que es lo que se argumentó en la
sección anterior, ahora recorrido en la dirección inversa---. $\vec P$ tiene las
mismas unidades que $\vec D$: $\si{\coulomb\per\metre\squared}$.
\end{notebox}

\subsection{Cargas de polarización}

Para la densidad volumétrica hay que usar la divergencia en \textbf{coordenadas
esféricas}, no la derivada radial a secas:
\[
\rho_p = -\div \vec P
= -\frac{1}{r^2}\pdv{r}\left(r^2 P_r\right).
\]

\begin{warnbox}[El $1/r^2$ que uno se come]
Es justamente el factor que se pierde si se escribe la fórmula de memoria. No hay
componentes según $\hat e_\theta$ ni $\hat e_\phi$, pero el $r^2$ sí está.
\end{warnbox}

Y ahí ocurre la cancelación: $P_r \propto 1/r^2$, así que $r^2 P_r$ es una
\textbf{constante} y su derivada se anula:
\[
\boxed{\ \rho_p = 0\ }
\]

\begin{notebox}[No era obvio]
Es la compensación de un $r^2$ con otro $r^2$, y no siempre tiene por qué pasar.
\end{notebox}

Para la densidad superficial, $\sigma_p = \vec P \cdot \hat n$ con $\hat n$
\textbf{saliente del material}. Acá está la sutileza del ejemplo: el material
dieléctrico está \textbf{afuera} de la esfera, de modo que la normal saliente del
material sobre la superficie $r=a$ apunta hacia el centro:
\[
\hat n = -\hat e_r
\quad\Longrightarrow\quad
\boxed{\ \sigma_p = -\frac{(K-1)\,Q}{4\pi K a^2}\ }
\]

\begin{warnbox}[Siempre hay dos normales]
La que va en la definición de $\sigma_p$ es la que \textbf{sale del material}, y
como acá el material está afuera, ese vector apunta hacia adentro. El signo menos
no es un descuido: es el contenido físico del resultado.
\end{warnbox}

Es \textbf{negativa} si $Q$ es positiva, y es razonable: frente a una carga
positiva, las moléculas se orientan de modo que quede su lado negativo hacia
ella. Todo el resto se compensa. La carga de polarización total, con densidad
uniforme sobre la esfera de radio $a$:
\[
Q_p = \sigma_p \cdot 4\pi a^2 = -\frac{K-1}{K}\,Q
= -\left(1 - \frac{1}{K}\right) Q .
\]

\subsection{Por qué el campo es menor}

Con esto se cierra el argumento físico que quedó pendiente. La carga
\textbf{total} que ve el campo eléctrico no es $Q$ sino

\begin{keybox}[La carga que se ve]
\[
Q + Q_p = Q - \left(1-\frac{1}{K}\right) Q = \frac{Q}{K},
\]
la carga original dividida por $K$ --- exactamente el factor por el que se redujo
$E$.
\end{keybox}

La imagen es la caricatura de siempre: la carga $+Q$ queda rodeada de cargas
negativas de polarización pegadas a ella, y la carga neta que se «ve» desde
afuera es menor.

\begin{notebox}[$\vec D$ no cuenta la polarización, $\vec E$ sí]
Para $\vec D$ la fuente es sólo la carga libre; para $\vec E$ hay que poner la
carga total, que es menor. De ahí que $D$ no lleve $K$ y $E$ sí.
\end{notebox}

Es un ejemplo de juguete ---podría haberse hecho en Física III--- pero fija las
ideas antes de pasar a los problemas propios de este curso.

\section{Condiciones de borde entre dos medios}

Los ejemplos que vienen usan la ecuación de Laplace, y ahí aparece un problema
nuevo: cuando hay \textbf{dos medios} ---dos dieléctricos, o un conductor y un
dieléctrico--- la ecuación no vale en la frontera. Hay que saber cómo se pegan
las soluciones de un lado y del otro.

\begin{figure}[H]
  \centering
  % Clase 10 — las dos construcciones de las condiciones de borde, lado a lado.
% Van juntas a proposito: son dos superficies distintas (una cerrada, una curva)
% sobre la MISMA interfaz, y lo que las distingue es que una controla la
% componente normal y la otra la tangencial. Separarlas pierde el contraste.
% Las construcciones se dibujan MAS GRANDES que la version anterior: cuando la
% capsula media 1,2 x 0,8 los rotulos S1/S2/lateral no entraban y se montaban
% sobre la interfaz. Con mas aire cada uno tiene su lugar.
% La prosa explicativa va en el \caption; adentro solo las dos conclusiones.
\begin{tikzpicture}[scale=1]

  % ---------- panel izquierdo: capsula gaussiana -> D_n ----------
  \begin{scope}[shift={(0,0)}]
    \fill[figblue!8] (-2.5,0) .. controls (-1.1,0.34) and (1.1,-0.34) .. (2.5,0)
      -- (2.5,2.1) -- (-2.5,2.1) -- cycle;
    \fill[figamber!12] (-2.5,0) .. controls (-1.1,0.34) and (1.1,-0.34) .. (2.5,0)
      -- (2.5,-2.1) -- (-2.5,-2.1) -- cycle;
    \draw[figline, line width=1.0pt]
      (-2.5,0) .. controls (-1.1,0.34) and (1.1,-0.34) .. (2.5,0);
    \node[figlbl, figblue] at (-1.75,1.72) {medio 1};
    \node[figlbl, figamber] at (-1.75,-1.72) {medio 2};

    % capsula, con aire suficiente para los rotulos
    \draw[figred, line width=1.0pt] (-0.85,-0.62) rectangle (0.85,0.55);
    \node[figlblaux, figred, fill=figblue!8, inner sep=1.5pt] at (-1.35,0.55) {$S_1$};
    \draw[figaux] (-1.05,0.55) -- (-0.6,0.55);
    \node[figlblaux, figred, fill=figamber!12, inner sep=1.5pt] at (-1.35,-0.62) {$S_2$};
    \draw[figaux] (-1.05,-0.62) -- (-0.6,-0.62);
    \node[figlblaux, figred] at (1.78,0.55) {lateral};
    \draw[figaux] (1.4,0.48) -- (0.9,0.1);

    \draw[figvecres] (-0.42,0.55) -- (-0.42,1.18);
    \node[figlbl, figred] at (-0.75,1.02) {$\hat n$};

    \node[figlbl, align=center] at (0,-2.75) {Gauss para diel\'ectricos};
    \node[figlbl, align=center] at (0,-3.4)
      {$D_{1n}-D_{2n}=\sigma_{\text{libre}}$};
  \end{scope}

  % ---------- panel derecho: curva cerrada -> E_t ----------
  \begin{scope}[shift={(6.8,0)}]
    \fill[figblue!8] (-2.5,0) .. controls (-1.1,0.34) and (1.1,-0.34) .. (2.5,0)
      -- (2.5,2.1) -- (-2.5,2.1) -- cycle;
    \fill[figamber!12] (-2.5,0) .. controls (-1.1,0.34) and (1.1,-0.34) .. (2.5,0)
      -- (2.5,-2.1) -- (-2.5,-2.1) -- cycle;
    \draw[figline, line width=1.0pt]
      (-2.5,0) .. controls (-1.1,0.34) and (1.1,-0.34) .. (2.5,0);
    \node[figlbl, figblue] at (-1.75,1.72) {medio 1};
    \node[figlbl, figamber] at (-1.75,-1.72) {medio 2};

    % rectangulo de altura pequena
    \draw[figvecres] (-0.85,0.26) -- (0.85,0.2);
    \draw[figvecres] (0.85,0.2) -- (0.85,-0.28);
    \draw[figvecres] (0.85,-0.28) -- (-0.85,-0.22);
    \draw[figvecres] (-0.85,-0.22) -- (-0.85,0.26);
    \node[figlbl, figred, fill=figblue!8, inner sep=1.5pt] at (0,0.72) {$C$};
    \node[figlblaux, figred] at (1.85,0.62) {alto $\to 0$};
    \draw[figaux] (1.42,0.52) -- (0.95,0.12);

    \draw[figvecaux] (-0.85,-0.85) -- (0.85,-0.85);
    \node[figlblaux, figgray] at (0,-1.18) {$\hat l$};

    \node[figlbl, align=center] at (0,-2.75) {$\oint_C \vec E\cdot d\vec l=0$};
    \node[figlbl, align=center] at (0,-3.4) {$E_{1t}=E_{2t}$};
  \end{scope}

\end{tikzpicture}

  \caption{Dos construcciones sobre la misma interfaz: una superficie cerrada
  controla la componente normal, una curva cerrada la tangencial. $\vec E$ sí
  puede saltar al cruzar, aun con $\sigma_{\text{libre}}=0$, si $K_1 \neq K_2$;
  pero el salto es finito, y una función con derivada acotada es continua, que
  es por lo que $\phi_1 = \phi_2$ puede reemplazar a la condición sobre $E_t$.}
\end{figure}

\subsection{La componente normal de \texorpdfstring{$\vec D$}{D}}

Se construye una \textbf{cápsula gaussiana} (un cilindro chato a caballo de la
interfaz) con tapa $S_1$ en el medio 1, tapa $S_2$ en el medio 2 y superficie
lateral. Tomándola suficientemente pequeña, la superficie se puede reemplazar por
un plano y $\vec D$ por una constante sobre cada tapa. La lateral es despreciable
frente a las tapas.

Aplicando Gauss para dieléctricos, con $\hat n$ el vector que va del medio 2 al 1:
\[
D_{1n}\,\Delta S - D_{2n}\,\Delta S = \sigma_{\text{libre}}\,\Delta S
\quad\Longrightarrow\quad
\boxed{\ D_{1n} - D_{2n} = \sigma_{\text{libre}}\ }
\]

\begin{warnbox}[La convención de signos importa]
Se puede elegir cualquiera de las dos normales, pero hay que ser consistente o
los signos cambian.
\end{warnbox}

Si no hay densidad superficial de carga \textbf{libre}, la componente normal de
$\vec D$ es continua.

\subsection{La componente tangencial de \texorpdfstring{$\vec E$}{E}}

Con la componente normal no alcanza: hace falta también la tangencial. Se toma un
\textbf{rectángulo pequeño} $C$ a caballo de la interfaz y se usa que la
circulación del campo eléctrico sobre cualquier curva cerrada es nula. Haciendo
tender a cero la altura del rectángulo:
\[
\boxed{\ E_{1t} = E_{2t}\ }
\]

En rigor no es una sola condición sino varias: vale para cualquier vector
tangencial, y a priori hay dos direcciones tangenciales independientes.

\subsection{La continuidad del potencial}

La condición tangencial admite un reemplazo más cómodo. Las dos condiciones
anteriores dicen que $\vec E$ \textbf{puede saltar} al cruzar la interfaz: aun
con $\sigma_{\text{libre}} = 0$, si las constantes dieléctricas difieren, la
continuidad de $D_n$ obliga a $E_n$ a dar un salto.

Pero lo importante es que esas discontinuidades son \textbf{finitas}. Y una
función cuya derivada tiene discontinuidades finitas es \textbf{continua}: se la
puede ver como la primitiva de esa derivada, y la primitiva es continua aunque no
derivable ---tendrá un pico, nada más---.

\begin{keybox}[Continuidad del potencial]
\[
\boxed{\ \phi_1 = \phi_2 \ \text{en la interfaz}\ }
\]
\end{keybox}

\begin{warnbox}[Es esencial que sea una superficie]
Cerca de una carga puntual el campo diverge; cerca de una línea de carga diverge
también, más suavemente, como un logaritmo. A través de una \textbf{superficie}
la discontinuidad es finita, y por eso el potencial es continuo.
\end{warnbox}

En la práctica se usa el par $\{\phi$ continuo, $D_n$ continua$\}$: la primera
condición es más fácil de imponer que la tangencial de $\vec E$, y la condición
sobre $D_n$ \textbf{no} se puede reemplazar.

\section{Poisson y Laplace en un dieléctrico}

En una zona ocupada por un medio homogéneo, isótropo y lineal de permitividad
$\varepsilon$, con $\vec D = \varepsilon \vec E$ y $\vec E = -\grad \phi$ (el
problema sigue siendo electrostático), la forma diferencial de Gauss da
$\div(\varepsilon \vec E) = \rho_{\text{libre}}$. Como $\varepsilon$ es
constante, sale de la divergencia:

\begin{keybox}[Poisson en un dieléctrico]
\[
\boxed{\ \laplacian \phi = -\frac{\rho_{\text{libre}}}{\varepsilon}\ }
\]
idéntica a la usual, con $\varepsilon$ donde antes había $\varepsilon_0$.
\end{keybox}

\begin{notebox}[Sólo carga libre]
De la carga de polarización no hay que ocuparse: ya está escondida en
$\varepsilon$.
\end{notebox}

En particular, si no hay carga libre ---aunque el medio \textbf{sí} esté
polarizado--- vale la ecuación de Laplace:
\[
\rho_{\text{libre}} = 0 \quad\Longrightarrow\quad \laplacian \phi = 0 .
\]

Aparece entonces una familia de problemas nueva: electrostática en presencia de
dieléctricos, donde sigue valiendo Laplace pero hay más de una región.

\section{Segundo ejemplo: cilindro en campo uniforme}

Un \textbf{cilindro muy largo} de radio $a$, de material dieléctrico de constante
$K$ y \textbf{descargado}, colocado en un campo eléctrico uniforme $\vec E_0$
perpendicular a su eje. Se pide el campo eléctrico en todo el espacio, adentro y
afuera.

\begin{notebox}[Dos campos, no uno]
El material se polariza, así que el campo deja de ser uniforme: tendrá la
contribución aplicada más el efecto de la polarización. Y como no es un
conductor, hay campo \textbf{adentro y afuera}. El ejemplo con esfera ya está
publicado en OpenFING de años anteriores; acá se hace con cilindro para variar,
con la misma estructura de razonamiento.
\end{notebox}

\begin{figure}[H]
  \centering
  % Clase 10 — el planteo del problema del cilindro.
% VERSION 2. La primera intentaba anclar las cinco condiciones A-E sobre el
% dibujo, y fue un error por dos motivos: (a) los cinco rotulos mas los de medio
% no entran sin pisarse, y (b) la tabla del texto ya las lista con su ubicacion,
% asi que era redundante. La figura se queda con lo que solo ella puede dar: la
% GEOMETRIA. Dos medios, dos potenciales, y las coordenadas.
% El docente aclara que "en el dibujo no se nota pero es un cilindro": de ahi la
% insinuacion del eje y la nota de que es muy largo.
\begin{tikzpicture}[scale=1]

  % zona exterior
  \fill[figblue!5] (-5.0,-2.7) rectangle (5.0,2.7);

  % campo aplicado
  \foreach \y in {-2.1,-1.05,0,1.05,2.1}{
    \draw[figvecaux] (-4.8,\y) -- (-3.5,\y);}
  \node[figlbl, figgray] at (-4.15,2.45) {$\vec E_0$};

  % cilindro visto de frente
  \fill[figamber!14] (0,0) circle (1.3);
  \draw[figamber, line width=1.0pt] (0,0) circle (1.3);

  % insinuacion del eje: es un cilindro largo, no un disco
  \draw[figaux] (-0.65,1.13) -- (0.35,1.82);
  \draw[figaux] (0.65,-1.13) -- (1.65,-0.44);
  \draw[figaux] (0.35,1.82) arc (115:-115:0.3 and 1.13);

  % medio exterior
  \node[figlbl, figblue] at (-2.75,1.15) {medio 1};
  \node[figlbl, figblue] at (-2.75,0.62) {$\varepsilon_0$};
  \node[figlbl, figblue] at (-2.75,0.05) {$\phi_1$};

  % medio interior
  \node[figlbl, figamber] at (0,-0.42) {$\phi_2$};
  \node[figlbl, figamber] at (3.15,-1.65) {medio 2:\ $\varepsilon$,\ $K$};
  \draw[figaux] (2.35,-1.5) -- (0.95,-0.95);

  % coordenadas, todas en el cuadrante superior derecho para no competir
  \draw[figaxis] (0,0) -- (3.3,0);
  \node[figlblaux] at (3.55,-0.05) {$x$};
  \draw[figvecaux] (0,0) -- (58:2.35);
  \node[figlblaux, fill=figblue!5, inner sep=1pt] at (58:1.95) {$r$};
  \draw[figgray, line width=0.5pt] (0.7,0) arc (0:58:0.7);
  \node[figlblaux] at (26:0.95) {$\theta$};

  % radio
  \draw[figvecaux] (0,0) -- (215:1.3);
  \node[figlblaux, fill=figamber!14, inner sep=1pt] at (215:0.75) {$a$};

  \node[figlblaux, align=center] at (2.7,2.15)
    {cilindro muy largo:\\ nada depende de $z$};

\end{tikzpicture}

  \caption{El planteo: dos zonas separadas por $r=a$, dos potenciales incógnita,
  y las coordenadas del problema. Las cinco condiciones de borde y dónde se
  aplica cada una están en la tabla de la sección siguiente.}
\end{figure}

\subsection{Las dos zonas}

Al ser un cilindro muy largo el problema no depende de $z$, y $\theta$ recorre de
$0$ a $2\pi$. La solución general de Laplace en cilíndricas bajo esas hipótesis
---los \textbf{armónicos cilíndricos}, vistos clases atrás--- es
\[
\phi(r,\theta) = a_0 + b_0 \ln r
+ \sum_{n=1}^{\infty} \left(a_n r^n + b_n r^{-n}\right)
\left(c_n \cos n\theta + d_n \sin n\theta\right).
\]

Como la permitividad cambia abruptamente en $r=a$, hacen falta \textbf{dos}
soluciones, con el mismo formato pero coeficientes distintos: $\phi_1$ afuera
(medio 1, $\varepsilon_0$) y $\phi_2$ adentro (medio 2, $\varepsilon$).

\subsection{Las cinco condiciones}

\begin{table}[H]
  \centering
  \begin{tabular}{clc}
    \toprule
     & Condición & Dónde \\
    \midrule
    \textbf{A} & $\phi_1 \sim -E_0\, r\cos\theta$      & $r \to \infty$ \\
    \textbf{B} & $b_0^{(1)} = 0$                       & $r \to \infty$ \\
    \textbf{C} & $\phi_2$ regular                      & $r \to 0$ \\
    \textbf{D} & $\phi_1(a,\theta) = \phi_2(a,\theta)$ & $r = a$ \\
    \textbf{E} & $D_{1n} = D_{2n}$                     & $r = a$ \\
    \bottomrule
  \end{tabular}
\end{table}

\textbf{A.} A gran distancia el campo debe tender al aplicado, cuyo potencial es
$-E_0 x = -E_0\, r\cos\theta$.

\textbf{B.} El término logarítmico es el potencial de una línea cargada, con
coeficiente proporcional a la carga por unidad de longitud. Como el cilindro está
descargado, ese coeficiente es nulo.

\begin{warnbox}[Por qué A no implica B]
Es un punto donde la intuición falla fácil. La condición A es un
\textbf{equivalente}, no un límite, y el logaritmo es \textbf{subdominante}
frente a $r$: crece mucho más despacio. Por eso A \textbf{no dice nada} sobre el
término logarítmico y hace falta B aparte. Lo que A sí mata son los términos que
crecen más rápido que $r$ ($r^2$, $r^3$, \dots).
\end{warnbox}

\begin{notebox}[Si el cilindro estuviera cargado]
B se reemplazaría por la condición de que a gran distancia se comporte como un
cilindro infinitamente fino: de lejos no se le ve el ancho.
\end{notebox}

\textbf{C.} En el origen no hay carga libre ni carga concentrada de ninguna
clase, y sólo una carga concentrada puede producir una divergencia. Por lo tanto
$\phi_2$ no diverge en $r \to 0$.

\textbf{D} y \textbf{E} son las condiciones de frontera de la sección anterior,
con $\sigma_{\text{libre}}=0$ porque el cilindro está descargado.

\begin{notebox}[La polarización superficial no cuenta acá]
Sí habrá densidad superficial \emph{de polarización}, pero esa no entra en la
condición sobre $\vec D$: ahí sólo cuenta la carga libre.
\end{notebox}

\subsection{Imposición y resultado}

\textbf{Con A y B}, en la zona exterior mueren $b_0^{(1)}$ y todos los
$a_n^{(1)}$ con $n \ge 2$. Del término lineal queda
$a_1^{(1)} r \left(c_1^{(1)}\cos\theta + d_1^{(1)}\sin\theta\right)
\sim -E_0 r\cos\theta$, que obliga a $d_1^{(1)} = 0$ y
$a_1^{(1)} c_1^{(1)} = -E_0$. Queda
\[
\phi_1(r,\theta) = a_0^{(1)} - E_0\, r\cos\theta
+ \sum_{n=1}^{\infty} \frac{b_n^{(1)}}{r^{n}}
\left(c_n^{(1)} \cos n\theta + d_n^{(1)} \sin n\theta\right).
\]

\textbf{Con C}, en la zona interior mueren $b_0^{(2)}$ y \textbf{todos} los
$b_n^{(2)}$ ---el logaritmo y todas las potencias negativas:
\[
\phi_2(r,\theta) = a_0^{(2)}
+ \sum_{n=1}^{\infty} a_n^{(2)} r^{n}
\left(c_n^{(2)} \cos n\theta + d_n^{(2)} \sin n\theta\right).
\]

\begin{notebox}[La constante que sí se puede elegir]
En cada zona se mató la mitad de los coeficientes. Además, el potencial está
definido a menos de una constante, así que se puede elegir $a_0^{(1)} = 0$ ---
pero sólo una de las dos: $a_0^{(2)}$ queda determinado por las condiciones.
\end{notebox}

\textbf{Con D}, se evalúan ambas en $r=a$ y se igualan para todo $\theta$. Como
$\{1, \cos n\theta, \sin n\theta\}$ son \textbf{linealmente independientes}, se
puede igualar coeficiente a coeficiente. Eso da $a_0^{(2)} = 0$, una ecuación por
cada seno, una por cada coseno con $n\ge 2$, y un caso aparte para $n=1$:
\[
-E_0\,a + \frac{b_1^{(1)} c_1^{(1)}}{a} = a_1^{(2)} c_1^{(2)}\, a .
\]

\begin{notebox}[Todavía no alcanza]
Son infinitas ecuaciones y sólo determinan las constantes de la zona 2 en
términos de las de la zona 1. Falta la condición E.
\end{notebox}

\textbf{Con E}, la componente normal es la radial, así que la condición es
$\varepsilon_0 \left.\pdv{\phi_1}{r}\right|_{r=a}
= \varepsilon \left.\pdv{\phi_2}{r}\right|_{r=a}$. Derivando ambas series e
igualando coeficiente a coeficiente aparece el mismo juego, ahora con los
factores $\varepsilon_0$ y $\varepsilon$ a cada lado.

Y ahí llegan las buenas noticias: combinando las ecuaciones de D y E, para cada
$n$ queda un sistema \textbf{sobredeterminado} cuya única solución es la trivial:
\[
d_n^{(1)} = d_n^{(2)} = 0 \ \ \forall n,
\qquad
c_n^{(1)} = c_n^{(2)} = 0 \ \ \forall n \ge 2 .
\]

Sobreviven únicamente los coeficientes de $n=1$, cuyas dos ecuaciones son
\[
-E_0\,a + \frac{b_1^{(1)} c_1^{(1)}}{a} = a\, a_1^{(2)} c_1^{(2)},
\qquad
-E_0 - \frac{b_1^{(1)} c_1^{(1)}}{a^2} = K\, a_1^{(2)} c_1^{(2)} .
\]

Sumando y restando convenientemente se despejan las dos incógnitas.

\begin{keybox}[Solución]
\[
\boxed{\ \phi_1(r,\theta) = -E_0\, r\cos\theta
+ a^2 E_0\,\frac{K-1}{K+1}\,\frac{\cos\theta}{r}\ }
\]
\[
\boxed{\ \phi_2(r,\theta) = -\frac{2 E_0}{K+1}\, r\cos\theta\ }
\]
\end{keybox}

\begin{figure}[H]
  \centering
  % Clase 10 — que dicen los dos potenciales que quedaron.
% CUIDADO CON EL ALCANCE: el docente cierra la clase diciendo que la discusion
% de la solucion (campo electrico, cargas) queda para la clase siguiente. Esta
% figura por lo tanto NO discute lineas de campo ni cargas inducidas: se limita
% a leer la estructura de los dos phi que el si escribio en el pizarron, que es
% un gradiente de distancia. La nota al pie deja constancia de la postergacion.
% Las flechas interiores van deliberadamente mas cortas que las exteriores: esa
% diferencia de longitud ES el contenido (el campo adentro esta reducido por
% 2/(K+1) < 1, porque K > 1).
\begin{tikzpicture}[scale=1]

  % ---------- panel izquierdo: phi_1 = aplicado + correccion ----------
  \begin{scope}[shift={(0,0)}]
    \node[figlbl] at (0,3.0) {afuera: $\phi_1$};
    \fill[figblue!5] (-2.4,-2.4) rectangle (2.4,2.4);
    \fill[figamber!14] (0,0) circle (0.95);
    \draw[figamber, line width=0.9pt] (0,0) circle (0.95);

    \foreach \y in {-1.9,-0.95,0.95,1.9}{
      \draw[figvecaux] (-2.3,\y) -- (-1.3,\y);}
    \node[figlblaux, figgray] at (-1.8,2.25) {$-E_0 r\cos\theta$};

    % el termino correctivo, dipolar en 2D
    \draw[figvec, figred] (1.25,0.62) -- (2.1,0.42);
    \draw[figvec, figred] (1.25,-0.62) -- (2.1,-0.42);
    \draw[figvec, figred] (1.35,0) -- (2.25,0);
    \node[figlblaux, figred, align=center] at (1.15,-1.72)
      {$a^2E_0\dfrac{K-1}{K+1}\dfrac{\cos\theta}{r}$};
    \node[figlblaux, figred] at (1.15,-2.25) {decae como $1/r$};
  \end{scope}

  % ---------- panel derecho: phi_2 uniforme y reducido ----------
  \begin{scope}[shift={(6.2,0)}]
    \node[figlbl] at (0,3.0) {adentro: $\phi_2$};
    \fill[figblue!5] (-2.4,-2.4) rectangle (2.4,2.4);
    \fill[figamber!14] (0,0) circle (0.95);
    \draw[figamber, line width=0.9pt] (0,0) circle (0.95);

    % campo exterior, para comparar longitudes
    \foreach \y in {-1.9,1.9}{
      \draw[figvecaux] (-2.3,\y) -- (-1.0,\y);}
    \node[figlblaux, figgray] at (-1.65,2.25) {$E_0$};

    % campo interior: uniforme y MAS CORTO
    \foreach \y in {-0.45,0,0.45}{
      \draw[figvec, figred, line width=1.0pt] (-0.55,\y) -- (0.25,\y);}
    \node[figlblaux, figred, align=center] at (0.15,-1.72)
      {$\vec E_2=\dfrac{2E_0}{K+1}\,\hat x$};
    \node[figlblaux, figred] at (0.15,-2.25) {uniforme y menor};
  \end{scope}

  \node[figlbl, align=center] at (3.1,-3.3)
    {$\phi_2$ es lineal en $x$: el campo interior es \textbf{uniforme}.};
  \node[figlbl, align=center] at (3.1,-3.8)
    {Y como $K>1$, el factor $2/(K+1)$ es menor que uno:};
  \node[figlbl, align=center] at (3.1,-4.3)
    {el diel\'ectrico \emph{apantalla} el campo aplicado.};

\end{tikzpicture}

  \caption{Qué dice cada uno de los dos potenciales.}
\end{figure}

\begin{notebox}[El método, más allá de las cuentas]
Hubo que hacer muchas cuentas, pero en el principio es el mismo tipo de problema
ya resuelto antes. La única diferencia es que hay dos zonas: se resuelve en cada
una y se imponen los bordes. Y en el borde se imponen normalmente dos
condiciones, una de continuidad y otra sobre una derivada.
\end{notebox}

\begin{notebox}[Por qué no hizo falta imponer $E_t$ continua]
Esa condición se escribe como la derivada respecto de $\theta$, y una vez
impuesta la continuidad de $\phi$ para todo $\theta$, la continuidad de
$\partial\phi/\partial\theta$ sale automáticamente. Las dos que se impusieron
---continuidad y derivada radial--- ya dan todo.
\end{notebox}

\vspace{1em}
\noindent\emph{La discusión de esta solución ---cómo queda el campo eléctrico y
cómo quedan las cargas--- queda explícitamente postergada para la clase
siguiente.}

\end{document}
